\clearpage

\section{矩阵的相抵与相似}


\subsection{基本概念提要}

矩阵的\textbf{相抵}与\textbf{相似}
是矩阵的两种重要的等价关系.

矩阵$A_{m \times n}$与$B_{m \times n}$相抵，
当且仅当存在可逆矩阵$P_{m \times m}$和$Q_{n \times n}$，
使得$A = PBQ$.
由于左乘可逆矩阵$P$相当于进行一系列初等行变换，
由于右乘可逆矩阵$Q$相当于进行一系列初等列变换，
“矩阵$A$与$B$相抵”也可以描述为“$A$与$B$可以通过初等变化互相转化”，
即“矩阵$A$与$B$的秩相同”.

方阵$A_{n \times n}$与$B_{m \times m}$相似，
当且仅当存在可逆矩阵$P_{n \times n}$，
使得$A = P^{-1}BP$.
对于给定的方阵$A$，
如果存在可逆矩阵$P$,
使得$A$相似变换成一个对角矩阵$D = P^{-1}AP$，
那么称$A$是可对角化的.
此时$D$也称为$A$的相似标准型.

寻求相似标准型的过程
引向了特征值和特征向量的概念.
若$A$可对角化：

\[
D = 
\mathrm{diag}(\lambda_1 , \cdots , \lambda_n) 
= P^{-1}AP
~~~
\text{即}
~~~
PD = AP
\]


则按照矩阵乘法的规则，
$P$的第$i$个列向量$\bm{X}_i$满足方程$\lambda_i \bm{X}_i = A \bm{X}_i$.
从直观上，
这个方程表明\textbf{$A$对向量$\bm{X}_i$的作用仅仅是伸缩，
而不改变其方向，$\lambda_i$是伸缩比例.}
这样的数$\lambda_i$和向量$\bm{X}_i$分别称为$A$的\textbf{特征值}和\textbf{特征向量}.
对$A$做相似对角化也就是求解$A$的所有特征值和特征向量.
让我们简单地回顾计算流程：

\medskip

(1)~~求解特征多项式$|\lambda I - A|$的所有根$\lambda_i$，
这些根即是特征值.

(2)~~对每个特征值$\lambda_i$，
解线性方程组$(\lambda_i I_n - A)\bm{X} = \bm{0}$.
解空间即是属于$\lambda_i$的\textbf{特征子空间}.

(3)~~判断$A$是否可对角化：
检查特征多项式中$\lambda_i$这个根的重数（称为\textbf{代数重数}），
以及属于$\lambda_i$的特征子空间的维数（称为\textbf{几何重数}）.
\textbf{$A$可相似对角化当且仅当所有特征值$\lambda_i$的几何重数等于代数重数.}

(4)~~若$A$可对角化，
取每个特征子空间的一组基，
合并成$\mathbb{K}^n$的一组基（$n$个线性无关的特征向量），
它们以列向量组的形式拼成矩阵$P$.

\medskip

相抵和相似都是在研究同一个线性映射在不同基下的表示矩阵.
在线性代数中，
我们研究矩阵的性质，
其实是为了线性映射的性质.
但是，
同一个线性映射在不同的基下可以表示成不同的矩阵.
从这个角度来看，
矩阵只是一种表象.
在本章中，
我们要通过研究表象之间的转换，
来破除表象，
逐渐回归线性映射的本质.


\subsection{相似不变量}

矩阵的相似是同一线性变换的不同表示矩阵之间的关系.
因此，
\textbf{相似不变量实质上是指线性变换反映在矩阵上的性质，
或者干脆认为是线性变换本身的性质.}

在前几章中，我们学习了矩阵的\textbf{秩}和\textbf{行列式}.
这两个概念之所以重要，
就在于它们都是相似不变量：
秩是线性变换的像空间维数，
行列式是线性变换的有向体积缩放因子.
这样的表述其实就反映了秩和行列式是属于线性变换本身的性质，
源于矩阵而高于矩阵.

证明秩和行列式是相似不变量并不困难：

\medskip

(1)~矩阵的相似是一种特殊的相抵关系，
因此秩也是相似不变量.

(2)~若$A_{n \times n} \sim B_{n \times n}$，
则存在可逆矩阵$P$使得$A = P^{-1} B P$.
两边取行列式得到$|A| = |P^{-1}| |B| |P|$.
注意$|P^{-1}| = |P|^{-1}$，
因此$|A| = |B|$，
即行列式是相似不变量.

\medskip

当然，
还有很多别的相似不变量.
但这里我们先指出，
特征多项式是一个重要的相似不变量.
它来源于相似对角化的计算过程.
以下给出证明：

\medskip

(*)~若$A \sim B$，
则存在可逆矩阵$P$使得$A = P^{-1} B P$.
我们有：

\begin{eq1}
 |\lambda I - A|
    = |P^{-1} (\lambda I) P - P^{-1}BP|
    = |P^{-1} (\lambda I - B) P|
    = |P^{-1}| |\lambda I - B| |P|
    = |\lambda I - B| \\[12pt]
\end{eq1}

因此\textbf{特征多项式}是一个相似不变量.
“多项式”形式上看起来抽象，
但是蕴含的信息量很大.
许多相似不变量都可以由它衍生出来.
比如，
\textbf{矩阵的秩其实就是特征值$0$的代数重数}；
\textbf{特征多项式的常数项是$(-1)^n |A|$}.
更一般地，
从特征多项式可以知道，
\textbf{特征多项式的根（也就是特征值）及其代数重数都是相似不变量，
特征多项式的各项系数也都是相似不变量.}

另一个特例是矩阵的\textbf{迹}，
它是指矩阵的所有对角线元素之和.
实际上这就是特征多项式的\textbf{一次项系数}的相反数.
如果$A$有$n$个特征值，
那么矩阵的迹也可以看成$n$个特征值之和.

请注意，
虽然特征多项式包含的信息量很大，
但是\textbf{两个矩阵的特征多项式相同远远不是两个矩阵相似的充分条件}.
在线性代数的学习范围内，
重点要知道的例子是，
特征子空间的维数（几何重数）也是相似不变量，
就不在特征多项式的管辖范围之内.
比如以下两个矩阵：

\[
    \begin{pmatrix}
        1 & 0 \\[2pt]
        0 & 1
    \end{pmatrix}
    ~~~
    \begin{pmatrix}
        1 & 1 \\[2pt]
        0 & 1
    \end{pmatrix} \\[12pt]
\]

它们有共同的特征多项式$(\lambda - 1)^2$，
但它们并不相似.
事实上后一个矩阵是不可对角化的.


\subsection{矩阵可对角化的一般条件}

本节的主要任务是提出几条矩阵可对角化的\textbf{充要条件}.

首先要提醒的是，
矩阵可对角化有一条很显然但容易被忽略的必要条件：
\textbf{矩阵必须存在$n$个特征值}，
换言之，
\textbf{特征多项式必须有$n$个根（计重数）}.
对这一条件的检查是判断矩阵是否可对角化的第一步.
而且，
由于多项式根的数量与数域有关，
因此在这一步，
\textbf{必须先明确讨论的是哪个数域.}
比如，
考察以下的矩阵：

\[
\begin{pmatrix}
    1 & -1 \\[2pt]
    1 & 1 
\end{pmatrix}  \\[12pt]
\]

它可以对角化吗？
正确的回答是不知道，
因为这个矩阵既可以理解为$\mathbb{R}$上的矩阵，
也可以理解为$\mathbb{C}$上矩阵，
这会导致不一样的回答.
可以计算出这个矩阵的特征多项式是$\lambda^2 + 2\lambda + 2$.
它在$\mathbb{R}$中没有根，
但在$\mathbb{C}$中有$2$个根.
结果，
这个矩阵在$\mathbb{R}$中不可对角化，
但在$\mathbb{C}$中可对角化.

根据代数基本定理，
一个$n$次多项式总有$n$个复根.
因此复矩阵必定有$n$个特征值（计重数）.
对于实矩阵（或者一般数域$\mathbb{K}$上的矩阵）则没有这样的结论.

\medskip

考虑矩阵对角化的意义.
对角矩阵代表的是一类最简单的线性变换——\textbf{伸缩变换}.
一个矩阵若能对角化，
说明它实质上代表的就是最简单的伸缩变换，
只不过选取的基不太合适，
使得表示矩阵比较复杂，
一眼看不出是个伸缩变换.
若能找出$n$个特征向量构成一组基，
那么这个线性变换在这组基下的表示矩阵就是对角矩阵了.
根据这一思想，
下述的定理是不证自明的：

\medskip

\textbf{定理6.3.1}~~
\textit{$n$阶矩阵$A$可对角化的充要条件是$A$有$n$个线性无关的特征向量.}

\medskip

直接从定义也很好理解这一点.
之前讨论过，
在$A$的对角化$D = P^{-1}AP$中，
$P$的列向量组实质上就是$n$个特征向量.
我们要求$P$是可逆矩阵，
就相当于说这$n$个特征向量是线性无关的.

\bigskip

我们要把这条“废话”定理转换成一个切实可行的判断准则——
在计算过程中，
如何判断总共有几个线性无关的特征向量呢？
其实我们是先计算特征子空间，
即$(\lambda I - A)\bm{X} = \bm{0}$的解空间，
它的一组基就是一组线性无关的特征向量.
这样一来，
每个不同的特征值$\lambda_i$都能提供$n_i$个线性无关的特征向量，
其中$n_i$是特征子空间$V_{\lambda_i}$的维数.
如果$A$有$m$个不同的特征值$\lambda_1 , \cdots , \lambda_m$，
那么总共会提供$n_1 + n_2 + \cdots + n_m$个线性无关的特征向量
——想一想这是为什么？
请注意一个简单的事实：
\textbf{属于不同特征值的特征子空间交集为$\{ \bm{0} \}$}
（一个特征向量不可能同时属于两个不同的特征值），
不同的特征子空间的和是直和.
从这个角度给出证明是非常简单的：

\medskip

\textbf{定理6.3.2}~~
\textit{设$\{ \bm{\alpha}_1 , \cdots , \bm{\alpha}_s \}$
和$\{ \bm{\beta}_1 , \cdots , \bm{\beta}_r \}$分别是矩阵$A$的属于$\lambda_1, \lambda_2$的线性无关的特征向量，
则合并之后的向量组$\{ \bm{\alpha}_1 , \cdots , \bm{\alpha}_s,\bm{\beta}_1 , \cdots , \bm{\beta}_r  \}$也线性无关.}

\textbf{证明\  }
由两个向量组的线性无关性可知
$\dim \langle \bm{\alpha}_1 , \cdots , \bm{\alpha}_s \rangle = s, \dim \langle \bm{\beta}_1 , \cdots , \bm{\beta}_r \rangle = r$.
则有：

\[
\dim \langle  \bm{\alpha}_1 , \cdots , \bm{\alpha}_s, \bm{\beta}_1 , \cdots , \bm{\beta}_r \rangle
= \dim \langle \bm{\alpha}_1 , \cdots , \bm{\alpha}_s \rangle + \dim \langle \bm{\beta}_1 , \cdots , \bm{\beta}_r \rangle 
= s + r
\]

这里第一个等号使用了线性子空间直和的维数公式.
这个式子直接表明合并之后的向量组是线性无关的.
\qed

\medskip


我们据此改进定理6.3.1的表述，
就得到了以下定理：

\medskip

\textbf{定理6.3.3}~~
\textit{$n$阶矩阵$A$可对角化的充要条件是$A$的属于不同特征值的特征子空间维数之和等于$n$.}

\medskip

以上总结了教材上的叙述的定理，
我们在此基础上略微多做一些讨论.
我们把特征值$\lambda_i$作为特征多项式的根的重数称为\textbf{代数重数}，
把相应的特征子空间维数称为\textbf{几何重数}.
探究矩阵可对角化的问题可以从代数重数和几何重数的关系入手，
这个想法是源自于我们具体计算矩阵对角化的过程.
就“对角化”的目标而言，
代数重数可以说是我们“预期”$\lambda_i$这个特征值能提供的线性无关特征向量的个数，
而几何重数就是实际上得到的个数.
本节开头讲过，
所有代数重数之和为$n$是矩阵可对角化的一个最基本的前提条件.
在此条件之下，
“可对角化”其实就是一个
“现实符合预期”的故事.
一般情况下，
几何重数有可能小于代数重数，
但不可能大于代数重数，
即“现实有可能不符合预期，
但决不会超出预期”.
这里简述一下理由：
设$n$阶方阵$A$有特征值$\lambda_i$，
其几何重数为$r$，
于是取到$\lambda_i$的$r$个线性无关的特征向量
$\bm{X}_1 , \cdots , \bm{X}_r$.
将其扩充成$\mathbb{K}^n$的一组基$\bm{X}_1 , \cdots , \bm{X}_r , \bm{X}_{r+1} , \cdots , \bm{X}_n$.
并将其拼成一个可逆矩阵$P$.
这可以把$A$相似变换成一个分块上三角矩阵：

\[
P^{-1}AP = 
\begin{pmatrix}
    \lambda_i I_r & C \\[2pt]
    O & D 
\end{pmatrix} \\[12pt]
\]

显然其特征多项式中$\lambda_i$是至少$r$重根，
即$A$的代数重数至少是$r$.

\medskip

\textbf{定理6.3.4}~~
\textit{任意特征值的几何重数不超过代数重数}

\medskip

结合定理6.3.3，
我们就知道：

\medskip

\textbf{定理6.3.5}~~
\textit{数域$\mathbb{K}$上的$n$阶矩阵$A$可对角化的充要条件是：$A$在$\mathbb{K}$中有$n$个特征值，
且每个特征值的代数重数等于几何重数.}

\medskip

根据这个定理，
只要出现某个特征值的几何重数小于代数重数，
就可以直接判断该矩阵不可对角化了.

灵活利用定理6.3.2及其推论，
如果我们发现矩阵$A$有$n$个不同的特征值，
那么可以直接判断它是可对角化的.

这里我们讨论的都是矩阵可对角化的充要条件.



\subsection{实对称矩阵的性质}

\textbf{定理6.4.1}~~
\textit{实对称矩阵的每个复特征值都是实数.}

\textbf{推论}~~
\textit{实对称矩阵必有$n$个实特征值（计重数）}

\medskip

\textbf{定理6.4.2}~~
\textit{实对称矩阵的属于不同特征值的特征向量是正交的.}

\textbf{定理6.4.3}~~
\textit{实对称矩阵一定正交相似于对角矩阵.}



\subsection{习题}


\subsubsection{相抵标准型的计算}

\begin{example}
    求下列矩阵的相抵标准型：
    \[
    \begin{pmatrix}
        1 & -1 & -3 & 1 \\
        1 & -1 & 2 & -1 \\
        4 & -4 & 3 & -2 
    \end{pmatrix}
    \]
\end{example}

\begin{remark}
    求秩即可.
\end{remark}

\begin{solution}
    做初等行变换：
  \[
    \begin{pmatrix}
        1 & -1 & -3 & 1 \\
        1 & -1 & 2 & -1 \\
        4 & -4 & 3 & -2 
    \end{pmatrix}
    \to
    \begin{pmatrix}
        1 & -1 & -3 & 1 \\
        0 & 0 & 5 & -2 \\
        0 & 0 & 15 & -6 
    \end{pmatrix}
    \to
     \begin{pmatrix}
        1 & -1 & -3 & 1 \\
        0 & 0 & 5 & -2 \\
        0 & 0 & 0 & 0
    \end{pmatrix}
    \]

    该矩阵的秩为$2$.
    因此相抵标准型是：
    \[
     \begin{pmatrix}
        1 & 0 & 0 & 0 \\
        0 & 1 & 0 & 0 \\
        0 & 0 & 0 & 0 
    \end{pmatrix}
    \]
\end{solution}

\bigskip

\begin{example}
    判断下列两个矩阵是否相抵？
    \[
    A = 
    \begin{pmatrix}
        1 & -1 & -3 & 1 \\
        1 & -1 & 2 & -1 \\
        4 & -4 & 3 & -2 
    \end{pmatrix}
    ~~~~~
    B = 
    \begin{pmatrix}
        1 & 3 & -2 & 0 \\
        3 & -2 & 0 & 1 \\
        4 & 1 & -2 & 1 
    \end{pmatrix}
    \]
\end{example}

\begin{solution}
    上一问中计算出矩阵$A$的秩为$2$.
    对矩阵$B$做初等行变换：
    \[
    \begin{pmatrix}
        1 & 3 & -2 & 0 \\
        3 & -2 & 0 & 1 \\
        4 & 1 & -2 & 1 
    \end{pmatrix}
    \to
    \begin{pmatrix}
        1 & 3 & -2 & 0 \\
        0 & -11 & 6 & 1 \\
        0 & -11 & 6 & 1 
    \end{pmatrix}
    \to
    \begin{pmatrix}
        1 & 3 & -2 & 0 \\
        0 & -11 & 6 & 1 \\
        0 & 0 & 0 & 0 
    \end{pmatrix}
    \]

    $B$的秩也是$2$.
    因此$A$与$B$相抵.
\end{solution}


\bigskip

\begin{example}
    给定矩阵$A$：
    \[
    A = 
    \begin{pmatrix}
        1 & -1 & -3 & 1 \\
        1 & -1 & 2 & -1 \\
        4 & -4 & 3 & -2 
    \end{pmatrix} \\[6pt]
    \]

    设$A$的相抵标准型为$A_0$.
    求可逆矩阵$C_1,C_2$，
    使得$A_0 = C_1 A C_2$.
\end{example}
\begin{remark}
    等式$P = C_1 A C_2$可以理解为
    “对$A$做一系列行变换，再做一系列列变换，得到$P$”.
    这些行变换和列变换分别化身为矩阵$C_1$和$C_2$.
    因此，
    所要做的是用初等变换把$A$老老实实地一步步化成相抵标准型，
    但同时要利用单位矩阵把你所做的每一步操作“记录”下来.
    以求$C_1$为例，
    把一个$3$阶单位矩阵放在$A$的右侧，
    对$A$做初等行变换时，
    这个$3$阶矩阵也做相应的初等行变换，
    直至$A$化成了简化行阶梯形矩阵.
    此时在右侧得到的行变换后的$3$阶矩阵就是$C_1$.
\end{remark}

\begin{solution}
    先做初等行变换，通过右侧放置单位矩阵来记录所做的行变换，
    得到$C_1$：
    \begin{align*}
        \left ( \begin{array}{cccc:ccc}
        1 & -1 & -3 & 1 & 1 & 0 & 0  \\
        1 & -1 & 2 & -1 & 0 & 1 & 0  \\
        4 & -4 & 3 & -2 & 0 & 0 & 1
        \end{array} \right ) 
        \to
        \left ( \begin{array}{cccc:ccc}
            1 & -1 & -3 & 1 & 1 & 0 & 0  \\
            0 & 0 & 5 & -2 & -1 & 1 & 0  \\
            0 & 0 & 15 & -6 & -4 & 0 & 1
        \end{array} \right )  \\[12pt]
        \to
        \left ( \begin{array}{cccc:ccc}
            2 & -2 & -6 & 2 & 2 & 0 & 0  \\
            0 & 0 & 5 & -2 & -1 & 1 & 0  \\
            0 & 0 & 0 & 0 & -1 & -3 & 1
        \end{array} \right ) 
        \to 
        \left ( \begin{array}{cccc:ccc}
            2 & -2 & -1 & 0 & 1 & 1 & 0  \\
            0 & 0 & 5 & -2 & -1 & 1 & 0  \\
            0 & 0 & 0 & 0 & -1 & -3 & 1
        \end{array} \right ) 
    \end{align*}
    
    再做初等行变换，通过下侧放置单位矩阵来记录所做的列变换，
    得到$C_2$：
    \begin{align*}
        \left ( 
            \begin{array}{cccc}
            2 & -2 & -1 & 0  \\
            0 & 0 & 5 & -2    \\
            0 & 0 & 0 & 0  \\ \hdashline
            1 & 0 & 0 & 0 \\
            0 & 1 & 0 & 0 \\
            0 & 0 & 1 & 0 \\
            0 & 0 & 0 & 1 
        \end{array} \right ) 
        \to
        \left ( 
            \begin{array}{cccc}
            2 & 0 & 0 & 0  \\
            0 & 0 & 0 & -2    \\
            0 & 0 & 0 & 0  \\ \hdashline
            1 & 1 & 1/2 & 0 \\
            0 & 1 & 0 & 0 \\
            0 & 0 & 1 & 0 \\
            0 & 0 & 5/2 & 1 
        \end{array} \right ) 
        \to \left ( 
            \begin{array}{cccc}
            1 & 0 & 0 & 0  \\
            0 & 1 & 0 & 0    \\
            0 & 0 & 0 & 0  \\ \hdashline
            1/2 & 0 & 1/2 & 1 \\
            0 & 0 & 0 & 1 \\
            0 & 0 & 1 & 0 \\
            0 & -1/2 & 5/2 & 0 
        \end{array} \right ) 
    \end{align*}

    因此：
    \begin{align*}
        C_1 = 
        \begin{pmatrix}
            1 & 1 & 0  \\
            -1 & 1 & 0  \\
            -1 & -3 & 1
        \end{pmatrix}
        ~~~~~
        C_2 = 
        \begin{pmatrix}
            1/2 & 0 & 1/2 & 1 \\
            0 & 0 & 0 & 1 \\
            0 & 0 & 1 & 0 \\
            0 & -1/2 & 5/2 & 0 
        \end{pmatrix}
    \end{align*}
\end{solution}

\subsubsection{相抵标准型的应用}

\begin{example}
    （满秩分解定理）
    设$A_{m \times n}$的秩为$r$.
    证明：存在秩为$r$的矩阵$B_{m \times r}$和秩为$r$的矩阵$C_{r \times n}$，
    使得$A = BC$.
\end{example}
\begin{remark}
    如果$A$本身是相抵标准型，你能按题意做出这个分解吗？
\end{remark}

\begin{solution}
    存在可逆矩阵$P_{m \times m }$和$Q_{n \times n}$，
    使得：
    \begin{align*}
        PAQ  = 
        \begin{pmatrix}
            I_r & O \\
            O & O
        \end{pmatrix}
    \end{align*}

    而这个相抵标准型可以直接分解成秩为$r$的$m \times r$矩阵和$r \times n$矩阵相乘：
    \begin{align*}
        \begin{pmatrix}
            I_r & O \\
            O & O
        \end{pmatrix}
        = 
        \begin{pmatrix}
            I_r  \\
            O 
        \end{pmatrix}
        \begin{pmatrix}
            I_r & O
        \end{pmatrix}
    \end{align*}

    将上述两式整理成$A$的分解式:
    \begin{align*}
        A  = P^{-1}
        \begin{pmatrix}
            I_r  \\
            O 
        \end{pmatrix}
        \begin{pmatrix}
            I_r & O
        \end{pmatrix} Q^{-1}
    \end{align*}

    则题目所要求的矩阵$B$和$C$即为：
    \begin{align*}
        B = 
        P^{-1}
        \begin{pmatrix}
            I_r  \\
            O 
        \end{pmatrix}
        ~~~~~
        C = 
        \begin{pmatrix}
            I_r & O
        \end{pmatrix} Q^{-1}
    \end{align*}

\end{solution}


\bigskip
\bigskip


\begin{example}

    (1)~
    设$n$阶方阵$A$的秩为$r~(r<n)$.
    证明：存在秩为$n-r$的$n$阶非零矩阵$B$和$C$，
    使得$AB=O$，$CA=O$.

    (2)~
    设$n$阶方阵$A$的秩为$r~(r<n)$.
    证明：存在$n$阶非零方阵$B$，
    使得$AB=BA = O$.
\end{example}

\begin{proof}
    直接做第(2)问.
    存在可逆矩阵$P_{m \times m }$和$Q_{n \times n}$，
    使得：
    \begin{align*}
        PAQ  = 
        \begin{pmatrix}
            I_r & O \\
            O & O
        \end{pmatrix}
    \end{align*}

    对于这个相抵标准型，
    可以直接找到一个秩为$n-r$的矩阵$R = \mathrm{diag}(O,I_{n-r})$，
    左乘或右乘的结果都是零矩阵：
    \begin{align*}
        \begin{pmatrix}
            I_r & O \\
            O & O
        \end{pmatrix} 
        \begin{pmatrix}
            O & O \\
            O & I_{n-r}
        \end{pmatrix}
        = 
        \begin{pmatrix}
            O & O \\
            O & I_{n-r}
        \end{pmatrix}
        \begin{pmatrix}
            I_r & O \\
            O & O
        \end{pmatrix} 
        =
        O
    \end{align*}

    即：
    \begin{align*}
        PAQR = RPAQ = O
    \end{align*}

    因为$P$和$Q$可逆，
    上式推出：
    \begin{align*}
        A\boxed{QR} = \boxed{RP}A = O
    \end{align*}

    这样就足以解答第(1)问了.
    为了凑出第(2)问所要求的矩阵$B$，
    只需要在保持等式的条件下，
    把两个方框里的矩阵形式凑成一样的：
    \begin{align*}
        A\boxed{QRP} = \boxed{QRP}A = O
    \end{align*}

    则$QRP$就是题目所要求的矩阵$B$.
\end{proof}

\bigskip

\begin{example}
    设矩阵$A,B,C$分别是$m \times n , ~r \times s ,~ m \times s$矩阵.
    证明：矩阵方程$AX - YB = C$有解$X,Y$的充要条件是：
    \[
    \mathrm{rank}
    \begin{pmatrix}
        A & O \\
        O & B 
    \end{pmatrix}
    =
    \mathrm{rank}
    \begin{pmatrix}
        A & C \\
        O & B
    \end{pmatrix}
    \]
\end{example}

\begin{remark}
    两个矩阵的秩相等，
    意味着可以通过初等行列变换从左侧变换到右侧.
    如果能写出初等变换对应的矩阵，
    就等于是解出了方程.
\end{remark}

\begin{proof}
    
    (1)~必要性.
    已知存在$n \times s$矩阵$X$和$m \times r$矩阵$Y$，
    使得$AX -  YB = C$.
    这个等式告诉我们，
    可以做如下的初等变换：
    \begin{align*}
        \begin{pmatrix}
            A & O \\
            O & B 
        \end{pmatrix} \to
        \begin{pmatrix}
            A & AX \\
            O & B 
        \end{pmatrix} \to
        \begin{pmatrix}
            A & AX-YB \\
            O & B 
        \end{pmatrix} =
        \begin{pmatrix}
            A & C \\
            O & B 
        \end{pmatrix}
    \end{align*}
    这就证明了：
    \begin{align*}
        \mathrm{rank}
        \begin{pmatrix}
            A & O \\
            O & B 
        \end{pmatrix}=
        \mathrm{rank}
        \begin{pmatrix}
            A & C \\
            O & B 
        \end{pmatrix}
    \end{align*}

    \medskip

    (2)~充分性.
    已知上述的秩等式成立.
    我们尝试去求解初等行列变换，
    即求解$P$和$Q$，
    使得：
    \begin{align*}
        P
        \begin{pmatrix}
            A & O \\
            O & B 
        \end{pmatrix}
        Q =
        \begin{pmatrix}
            A & C \\
            O & B 
        \end{pmatrix}
    \end{align*}

    我们借助相抵标准型来完成这件事.
    设$A,B$的秩分别为$r,t$.
    则存在可逆矩阵$P_1 , Q_1 , P_2 , Q_2$，
    使得：
    \begin{align*}
        P_1 A Q_1 =
        \begin{pmatrix}
            I_r & O \\
            O & O 
        \end{pmatrix}
        ~~~~~
        P_2 B Q_2 =
        \begin{pmatrix}
            I_t & O \\
            O & O 
        \end{pmatrix}
    \end{align*}

    于是有：
    \begin{align*}
        \begin{pmatrix}
            P_1 & O \\
            O & P_2
        \end{pmatrix}
        \begin{pmatrix}
            A & O \\
            O & B
        \end{pmatrix}
        \begin{pmatrix}
            Q_1 & O \\
            O & Q_2
        \end{pmatrix}
        =
        \begin{pmatrix}
            I_r & O & O & O \\
            O & O & O & O \\
            O & O & I_t & O \\
            O & O & O & O \\
        \end{pmatrix}
    \end{align*}

    同样的变换施加到另一个矩阵上是：
    \begin{align*}
        \begin{pmatrix}
            P_1 & O \\
            O & P_2
        \end{pmatrix}
        \begin{pmatrix}
            A & C \\
            O & B
        \end{pmatrix}
        \begin{pmatrix}
            Q_1 & O \\
            O & Q_2
        \end{pmatrix}
        =
        \begin{pmatrix}
            P_1 A Q_1 & P_1 C Q_2 \\
            O & P_2 B Q_2
        \end{pmatrix}
        =
        \begin{pmatrix}
            I_r & O & C_1 & C_2 \\
            O & O & C_3 & C_4 \\
            O & O & I_t & O \\
            O & O & O & O \\
        \end{pmatrix}
    \end{align*}

    对它进一步施加初等行列变换，可以再简化：
    \begin{align*}
        \begin{pmatrix}
            I_r & O & C_1 & C_2 \\
            O & O & C_3 & C_4 \\
            O & O & I_t & O \\
            O & O & O & O \\
        \end{pmatrix}
        \to 
        \begin{pmatrix}
            I_r & O & O & C_2 \\
            O & O & O & C_4 \\
            O & O & I_t & O \\
            O & O & O & O \\
        \end{pmatrix}
        \to 
        \begin{pmatrix}
            I_r & O & O & O \\
            O & O & O & C_4 \\
            O & O & I_t & O \\
            O & O & O & O \\
        \end{pmatrix}
    \end{align*}

    也就是说存在矩阵$G_1,G_2$，
    使得：
    \begin{align*}
        \begin{pmatrix}
            I_m & G_1 \\
            O & I_r
        \end{pmatrix}
        \begin{pmatrix}
            I_r & O & C_1 & C_2 \\
            O & O & C_3 & C_4 \\
            O & O & I_t & O \\
            O & O & O & O \\
        \end{pmatrix}
        \begin{pmatrix}
            I_n & G_2 \\
            O & I_s
        \end{pmatrix}
        = 
        \begin{pmatrix}
            I_r & O & O & O \\
            O & O & O & C_4 \\
            O & O & I_t & O \\
            O & O & O & O \\
        \end{pmatrix}
    \end{align*}

    事实上可以直接读出矩阵$G_1,G_2$分别是：
    \begin{align*}
        G_1 = 
        \begin{pmatrix}
            -C_1 & O \\
            -C_2 & O
        \end{pmatrix}
        ~~~~~
        G_2 = 
        \begin{pmatrix}
            O & -C_2 \\
            O & O
        \end{pmatrix}
    \end{align*}

    从最右侧的形式可以看出，
    矩阵的秩为$r+t+\mathrm{rank}(C_4)$.
    但由已知条件知道它的秩应该是$r+t$.
    因此$C_4 = O$.
    这样我们就已经做到了将两个矩阵都通过初等行列变换变成$\mathrm{diag}(I_r,O,I_t,O)$，
    以此为桥梁来整理上述的所有式子：
    \begin{align*}
        &\mathrm{diag}(I_r,O,I_t,O)
        = \\[4pt]
        &
        \begin{pmatrix}
            P_1 & O \\
            O & P_2
        \end{pmatrix}
        \begin{pmatrix}
            A & O \\
            O & B
        \end{pmatrix}
        \begin{pmatrix}
            Q_1 & O \\
            O & Q_2
        \end{pmatrix}
        = 
        \begin{pmatrix}
            I_m & G_1 \\
            O & I_r
        \end{pmatrix}
         \begin{pmatrix}
            P_1 & O \\
            O & P_2
        \end{pmatrix}
        \begin{pmatrix}
            A & C \\
            O & B
        \end{pmatrix}
        \begin{pmatrix}
            Q_1 & O \\
            O & Q_2
        \end{pmatrix}
        \begin{pmatrix}
            I_n & G_2 \\
            O & I_s
        \end{pmatrix}
    \end{align*}

    整理并计算得到：
    \begin{align*}
        \begin{pmatrix}
            A & C \\
            O & B
        \end{pmatrix}
        &= 
        \begin{pmatrix}
            P_1^{-1} & O \\
            O & P_2^{-1}
        \end{pmatrix}
        \begin{pmatrix}
            I_m & G_1 \\
            O & I_r
        \end{pmatrix}^{-1}
        \begin{pmatrix}
            P_1 & O \\
            O & P_2
        \end{pmatrix}
        \begin{pmatrix}
            A & O \\
            O & B
        \end{pmatrix}
        \begin{pmatrix}
            Q_1 & O \\
            O & Q_2
        \end{pmatrix}
        \begin{pmatrix}
            I_n & G_2 \\
            O & I_s
        \end{pmatrix}^{-1}
        \begin{pmatrix}
            Q_1^{-1} & O \\
            O & Q_2^{-1}
        \end{pmatrix} \\[4pt]
        &= \begin{pmatrix}
            P_1^{-1} & O \\
            O & P_2^{-1}
        \end{pmatrix}
        \begin{pmatrix}
            I_m & -G_1 \\
            O & I_r
        \end{pmatrix}
        \begin{pmatrix}
            P_1 A Q_1 & O \\
            O & P_2 B Q_2
        \end{pmatrix}
        \begin{pmatrix}
            I_n & -G_2 \\
            O & I_s
        \end{pmatrix}
        \begin{pmatrix}
            Q_1^{-1} & O \\
            O & Q_2^{-1}
        \end{pmatrix} \\[4pt]
        &= \begin{pmatrix}
            P_1^{-1} & O \\
            O & P_2^{-1}
        \end{pmatrix}
        \begin{pmatrix}
            P_1 A Q_1 & -G_1 P_2 B Q_2 - P_1 A Q_1 G_2 \\
            O & P_2 B Q_2
        \end{pmatrix}
        \begin{pmatrix}
            Q_1^{-1} & O \\
            O & Q_2^{-1}
        \end{pmatrix} \\[4pt]
        &= 
        \begin{pmatrix}
            A & -P_1^{-1} G_1 P_2 B  - A Q_1 G_2 Q_2^{-1}\\
            O & B
        \end{pmatrix}
    \end{align*}

    比较左右两侧，得到：
    \begin{align*}
        C  =  -P_1^{-1} G_1 P_2 B  - A Q_1 G_2 Q_2^{-1}
    \end{align*}

    这样就找到了矩阵方程$AX - YB = C$的解：
    \begin{align*}
        X = -Q_1 G_2 Q_2^{-1},~~~Y= P_1^{-1} G_1 P_2
    \end{align*}

    因此矩阵方程是有解的.
\end{proof}


\subsubsection{相似对角化的计算}

\begin{example}
    设矩阵：
    \[
    A  =
    \begin{pmatrix}
        4 & 2 & 2 \\[2pt]
        0 & 4 & 0 \\[2pt]
        0 & -6 & -2 
    \end{pmatrix}  \\[12pt]
    \]

    (1)~求$A$的全部特征值与特征向量.

    (2)~判断$A$是否可对角化.
    如果可对角化，求可逆矩阵$U$和对角矩阵$D$，
    使得$D = U^{-1}AU$；
    如果不可对角化，说明理由.

    (3)~计算$A^{20}$.
\end{example}

\begin{solution}
    ~

    (1)计算特征多项式：
    \[
    |\lambda I - A|=
    \begin{vmatrix}
    \lambda-4 & -2 & -2\\
    0 & \lambda-4 & 0\\
    0 & 6 & \lambda+2
    \end{vmatrix}
    = (\lambda-4)^2 (\lambda + 2)
    \]

    $A$的特征值为：
    $\lambda_1=4$（$2$重）,
    $\lambda_2=-2$（$1$重）.

    (i) $\lambda=4$：
    解方程$(4I-A)\bm{X}=\bm{0}$.
    系数矩阵$4I-A$为：
    \begin{align*}
        \begin{pmatrix}
    0 & -2 & -2\\
    0 & 0 & 0\\
    0 & 6 & 6
    \end{pmatrix}
    \end{align*}

    它的秩为$1$，
    能得到$2$个线性无关的特征向量：
    \begin{align*}
        \bm{X}_1=\begin{pmatrix}1\\0\\0\end{pmatrix},\quad
    \bm{X}_2=\begin{pmatrix}0\\-1\\1\end{pmatrix}
    \end{align*}

    构成基础解系，
    即特征子空间的一组基.
    因此$\lambda = 4$对应的全部特征向量为$c_1 \bm{X}_1 + c_2 \bm{X}_2$，
    其中$c_1 ,c_2$不同时为$0$.


    (ii)~$\lambda=-2$： 
    解方程$(-2I-A)\bm{X}=\bm{0}$.
    系数矩阵$-2I-A$为：
    \begin{align*}
        \begin{pmatrix}
    6 & 2 & 2\\
    0 & 6 & 0\\
    0 & -6 & 0
    \end{pmatrix}
    \end{align*}

    它的秩为$2$，
    能得到$1$个线性无关的特征向量：
    \begin{align*}
        \bm{X}_3=\begin{pmatrix}1\\0\\-3\end{pmatrix}
    \end{align*}
    
    构成基础解系，
    即特征子空间的一组基.
    因此$\lambda = -2$对应的全部特征向量为$c_3 \bm{X}_3$，
    其中$c_3 \ne 0$.

    \medskip

    (2)~在(1)中，
    我们已经得到了$3$个线性无关的特征向量，
    因此$A$可对角化.

    \medskip

    (3)~
    设矩阵：
    \begin{align*}
         P= (\bm{X}_1 , \bm{X}_2 , \bm{X}_3) = 
        \begin{pmatrix} 
    1 & 0 & 1\\
    0 & -1 & 0\\
    0 & 1 & -3
    \end{pmatrix},
    \quad
    D=\begin{pmatrix}
    4 & 0 & 0\\
    0 & 4 & 0\\
    0 & 0 & -2
    \end{pmatrix}
    \end{align*}

    则有$D=P^{-1}AP$.
    于是：
    \begin{align*}
        A^{20} 
        = (PDP^{-1})^{20}
        = PD^{20}P
        = \begin{pmatrix}
        4^{20} & 2(4^{20}-2^{20}) & 2(4^{20}-2^{20})\\
        0 & 4^{20} & 0\\
        0 & -6(4^{20}-2^{20}) & 2^{20}
        \end{pmatrix}
    \end{align*}

\end{solution}

\bigskip

\begin{example}
    设实矩阵：
    \[
    A = 
    \begin{pmatrix}
        -4 & 2 & 2 \\[2pt]
        2 & -1 & 4 \\[2pt]
        2 & 4 & a
    \end{pmatrix} \\[12pt]
    \]

    已知$-5$是$A$的一个重数为$2$的特征值. 
    求$a$的值.
\end{example}

\begin{remark}
    固然可以直接去计算特征多项式.
    但是注意$A$是实对称矩阵，
    这应该能带来很大的方便.
\end{remark}

\begin{solution}
    注意$A$是实对称矩阵，
    $-5$是$A$的一个重数为$2$的特征值，
    意味着$-5$能对应$2$个线性无关的特征向量，
    即$-5I - A$的秩为$1$.
    计算$-5I - A$：
    \begin{align*}
        \begin{pmatrix}
            -1 & -2 & -2 \\
        -2 & -4 & -4 \\
        -2 & -4 & -5-a
        \end{pmatrix}
    \end{align*}

    由秩为$1$可知$a=-1$.
\end{solution}

\bigskip

\subsubsection{对角化相关证明}

\begin{example}
    设$n$阶矩阵：
    \begin{eq1}
        A = 
        \begin{pmatrix}
            I_r & B \\[2pt]
            O & -I_{n-r}
        \end{pmatrix}\\[12pt]
    \end{eq1}

    求证：$A$可对角化.
\end{example}

\begin{remark}
    我们所学的矩阵可对角化的理论直接来源于对角化的计算过程.
    对于本题这种形式具体的矩阵，
    直接去尝试计算它的对角化即可.
    目标是得到代数重数和几何重数的关系.
\end{remark}

\begin{proof}
    $A$是上三角矩阵，
    容易求得$A$的特征多项式为$|\lambda I - A| = (\lambda - 1)^r (\lambda + 1)^{n-r}$.
    可知$1$是$r$重特征值，
    $-1$是$n-r$重特征值.
    只需证明它们的特征子空间的维数（几何重数）也分别是$r$和$n-r$.

    所谓特征子空间的维数，
    就是齐次线性方程组$(\lambda I - A) \bm{X} = \bm{0}$的解空间维数.
    根据维数公式，这等于$n - \mathrm{rank}(\lambda I - A)$.
    于是我们考察矩阵$I - A$和$-I - A$的秩，
    易见它们分别是$n-r$和$r$.
    因此特征值$1$和$-1$的特征子空间维数分别为$r$和$n-r$.
\end{proof}

\bigskip

如果已知$\lambda_i$是特征值，
那么矩阵$A$的几何重数其实就是
$\dim \operatorname{Ker} (\lambda I - A)$，
根据维数公式，
几何重数还可以写成$n - \mathrm{rank}(\lambda I - A)$.
如果要证明矩阵$A$可对角化，
可以考虑利用充要条件，
去证所有特征值的几何重数之和为$n$即可.

明确了这个目标之后，
问题就转化成了证明一些矩阵的秩的关系.
具体处理办法见上一章.

\begin{example}
    ~

    (1)~~求证：若$n$阶方阵$A$是幂等矩阵（$A^2 = A$），则$A$必可对角化.

    (2)~~求证：若$n$阶方阵$A$是对合矩阵（$A^2 = I_n$），则$A$必可对角化.
\end{example}

\begin{proof}
    ~

    (1)~先计算特征值.
    设$\lambda$是$A$的特征值，
    对应的一个特征向量是$\bm{X}$.
    由$A^2 = A$可知$A^2 \bm{X} = A \bm{X} = \lambda \bm{X}$.
    但是另一方面，
    $A^2 \bm{X} = A (A \bm{X}) = A (\lambda \bm{X}) = \lambda A \bm{X} = \lambda^2 \bm{X}$.
    比较两式可知$\lambda = \lambda^2$，
    即$A$的特征值只可能是$0$或$1$.
    我们只考察它们的特征子空间维数之和是否为$n$，
    即证明$\dim \operatorname{Ker}(I - A) + \dim \operatorname{Ker} A = n$.

    根据维数公式：$n = \dim \operatorname{Ker}(I - A) + \dim \operatorname{Im}(I - A)$，
    $n = \dim \operatorname{Ker}A + \dim \operatorname{Im}A$.
    只需证明$\dim \operatorname{Ker}(I - A) + \dim \operatorname{Ker}A = n$，
    即$\mathrm{rank}(I-A) + \mathrm{rank}(A) = n$.
    这是上一章习题中证明过的结论.

    \medskip

    (2)~先计算特征值.
    设$\lambda$是$A$的特征值，
    对应的一个特征向量是$\bm{X}$.
    由$A^2 = I$可知$A^2 \bm{X} = I \bm{X} = \bm{X}$.
    但是另一方面，
    $A^2 \bm{X} = A (A \bm{X}) = A (\lambda \bm{X}) = \lambda A \bm{X} = \lambda^2 \bm{X}$.
    比较两式可知$\lambda^2 = 1$，
    即$A$的特征值只可能是$1$或$-1$.
    我们只考察它们的特征子空间维数之和是否为$n$，
    即证明$\dim \operatorname{Ker}(I - A) + \dim \operatorname{Ker} (I+A) = n$.

    根据维数公式：$n = \dim \operatorname{Ker}(I - A) + \dim \operatorname{Im}(I - A)$，
    $n = \dim \operatorname{Ker}(I+A) + \dim \operatorname{Im}(I+A)$.
    只需证明$\dim \operatorname{Ker}(I - A) + \dim \operatorname{Ker}(I+A) = n$，
    即$\mathrm{rank}(I-A) + \mathrm{rank}(I+A) = n$.
    这是上一章习题中证明过的结论.
\end{proof}

\bigskip

\begin{example}
    设$n$阶矩阵$A$和$B$满足$AB+B=O$，
    且有$\mathrm{rank}(B) = k~(1 \leq k \leq n-1)$.
    证明：$-1$是$A$的一个特征值，
    且$A$至少有$k$个属于$-1$的线性无关特征向量.
\end{example}

\begin{proof}
    我们的目标是研究$I+A$的秩.
    已知$A$和$B$满足$(A+I)B = O$.
    根据上一章习题中证明过的结论，
    $\mathrm{rank}(A+I) + \mathrm{rank}(B) \leq n$.
    再由维数公式：$\dim \operatorname{Ker} (A+I) + \mathrm{rank}(A+I) = n$.
    代入可知$\dim \operatorname{Ker} (A+I) \geq k$，
    即属于特征值$-1$的特征子空间至少是$k$维，
    亦即至少有$-1$个属于$-1$的线性无关的特征向量.
\end{proof}

\bigskip

\begin{example}
    设$\bm{\alpha},\bm{\beta}$是非零的$n$维列向量$(n \geq 2)$，
    且$A = \bm{\alpha} \bm{\beta}^T$.
    试问矩阵$A$是否相似于对角矩阵，
    给出理由.
\end{example}

\begin{remark}
    本题的切入点是注意形如$\bm{\alpha} \bm{\beta}^T$的矩阵秩为$1$.
\end{remark}

\begin{proof}
    按矩阵乘法的规则，
    $A = \bm{\alpha} \bm{\beta}^T$的每一个列向量都与$\bm{\alpha}$共线，
    因此$\mathrm{rank}(A) < 1$.
    又因为$\bm{\alpha},\bm{\beta}$都是非零向量，
    $A =  \bm{\alpha} \bm{\beta}^T$不是零矩阵.
    因此$\mathrm{rank} (A) = 1$.
    由维数公式可知$\dim \operatorname{Ker} A = n-1$，
    即$A$的$0$特征值的几何重数是$n-1$.
    这表明$0$的代数重数至少是$n-1$.

    设$f(\lambda)$是$A$的特征多项式.
    由上述讨论可知，
    $\lambda^{n-1}$是$f(\lambda)$的一个因式.
    而$f(\lambda)$是$n$次多项式.
    因此$f(\lambda)$是$\lambda^{n-1}$与某个一次因式$\lambda - a$的乘积.
    这就说明$f(\lambda)$一定还有一个根$a$，
    即$f(\lambda)$一定有$n$个特征值.

    对剩下的一个特征值$a$进行分类讨论：
    如果$a=0$，
    则特征值$0$的代数重数为$n$，
    而几何重数是$n-1$.
    则$A$不可对角化.
    此时$a \ne 0$，
    则特征值$0$的代数重数和几何重数都是$n-1$，
    而特征值$a$的代数重数和几何重数都是$1$.
    此时$A$可对角化.
\end{proof}


\bigskip

\begin{example}
    设$T$和$S$是数域$\mathbb{K}$上的$n$阶矩阵，
    已知$T$在$\mathbb{K}$上有$n$个互不相同的特征值.
    证明：若$ST = TS$，
    则$S$可对角化.
\end{example}

\begin{proof}
    $T$有$n$个互不相同的特征值，
    则$T$可对角化，
    且每个特征值的特征子空间都是$1$维的.
    设$\lambda$是$T$的一个特征值，
    对应的一个特征向量是$\bm{x}$.
    则有：
    \begin{align*}
        T (S \bm{x})
        = ST \bm{x} 
        = S (\lambda \bm{x})
        = \lambda (S \bm{x})
    \end{align*}

    可见$S\bm{x}$是$T$的属于$\lambda$的特征向量.
    因为$T$的属于$\lambda$的特征子空间是一维的，
    即$\langle \bm{x} \rangle$，
    因此存在数$\mu \in \mathbb{K}$，
    使得$S \bm{x} = \mu \bm{x}$.
    因此$\bm{x}$是$S$的特征向量.
    由$\lambda$的任意性可知，
    $T$特征向量都是$S$的特征向量，
    因此$S$有$n$个线性无关的特征向量，
    $S$可对角化.
\end{proof}

\bigskip

\begin{example}
    设$A,B$为$n$阶方阵. 
    求证：$AB-BA$必定不相似于$kI_n~(k \ne 0)$.
\end{example}

\begin{remark}
    只要两个矩阵的某个相似不变量不同，
    那它们就一定不相似.
\end{remark}

\begin{proof}
    $\mathrm{tr}(AB - BA) = \mathrm{tr}(AB) - \mathrm{tr}(BA) = 0$，
    但$\mathrm{tr}(kI_n) = nk$.
    $\mathrm{tr}(AB - BA) \ne \mathrm{tr}(kI_n)$，
    因此它们必定不相似.
\end{proof}

\subsubsection{*相似对角化的应用}

\begin{example}
    特征方程的概念在许多问题中都有应用.

    (1)~考虑下述矩阵$A$：
    \[
    \begin{pmatrix}
        2 & 3 \\
        1 & 0
    \end{pmatrix}
    \]

    求可逆矩阵$P$，
    使得$P^{-1}AP$是对角矩阵.

    (2)~数列$\{ a_n \}$满足递推公式：
    $a_{n} = 2a_{n-1} + 3a_{n-2}$.
    若已知前两项$a_1 , a_2$，
    利用(1)的结论计算通项公式.

    (3)~利用(1)的结论求解二阶常系数线性微分方程$y'' = 2y'+3y$.
\end{example}

\begin{solution}
    ~
    
    (1)~$A$的特征多项式为$\lambda^2 - 2 \lambda - 3 = 0$.
    特征值为$\lambda_1 = 3,~\lambda_2 = -1$.
    解得它们分别对应的一个特征向量是$\bm{\xi}_1 = (3,1)^T, ~\bm{\xi}_2 = (1,-1)^T$.
    因此题目所要求的矩阵$P$为：
    \begin{align*}
        \begin{pmatrix}
        3 & 1 \\[2pt]
        1 & -1
        \end{pmatrix}
    \end{align*}
    

    满足$P^{-1}AP = \mathrm{diag}(3,-1)$.

    \medskip

    (2)~记状态向量$\bm{\alpha}_n = (a_{n+1}, a_n)^T$.
    则递推公式可以表达为矩阵形式：
    \begin{align*}
        \begin{pmatrix}
            a_{n+2} \\
            a_{n+1}
        \end{pmatrix}
        =
        \begin{pmatrix}
        2 & 3 \\
        1 & 0
        \end{pmatrix}
        \begin{pmatrix}
            a_{n+1} \\
            a_{n}
        \end{pmatrix}
        ~~~~~\text{即}~
        \bm{\alpha}_{n+1}
        = 
        A\bm{\alpha}_1
    \end{align*}

\end{solution}